% =====================================================================
\section{Visual Processing Pipeline}
\label{sec:visual}
% =====================================================================

The visual pipeline converts a video clip into a temporally ordered sequence of
fixed-size mouth images and then into a sequence of visual feature vectors.
Preprocessing (\cref{sec:mouth-crop}) runs once per utterance and is cached;
the encoder (\cref{sec:visual-encoder}) is part of the trained model.

\subsection{Mouth extraction}
\label{sec:mouth-crop}

\paragraph{1. Reading frames.} Frames are decoded in order with
\texttt{cv2.VideoCapture} and converted from BGR to RGB, the format MediaPipe
expects. Each frame keeps its index $t$, which is the only time reference used
for GRID; at 25\,fps, frame $t$ starts at $40t$\,ms.

\paragraph{2. Detecting landmarks.} MediaPipe Face Mesh returns 468
three-dimensional landmarks per detected face~\cite{kartynnik2019,aripin2024}.
It is run in video mode with \texttt{max\_num\_faces=1}, so that landmarks are
tracked from frame to frame instead of being re-detected independently, which
is both faster and more temporally stable.

\paragraph{3. Selecting lip landmarks.} The Face Mesh topology defines a fixed
subset of landmarks on the outer and inner lip contours. The implementation
takes these indices from MediaPipe's own lip-connection definition
(\texttt{FACEMESH\_LIPS}) rather than hard-coding them, which avoids errors if
the index list is mistyped.

\paragraph{4. Computing the mouth box.} Let $(x_i, y_i)$ be the pixel
coordinates of the lip landmarks in frame $t$. The mouth centre and size are
\begin{equation}
c_t = \Big(\tfrac{\min_i x_i + \max_i x_i}{2},\; \tfrac{\min_i y_i + \max_i y_i}{2}\Big),
\qquad
S_t = \max\!\big(\max_i x_i - \min_i x_i,\; \max_i y_i - \min_i y_i\big).
\end{equation}
A square box is used so that resizing does not distort the mouth's aspect
ratio.

\paragraph{5. Adding a margin and stabilising.} The box side is enlarged to
$L_t = (1+2m)\,S_t$ with a margin $m = 0.25$ \rec, so that the chin, the corners
of the mouth and jaw movement remain visible when the mouth opens. Raw
per-frame boxes jitter by a few pixels, and this jitter would appear to the
network as spurious motion. The centre and side length are therefore smoothed
over time with a centred moving average of five frames \rec. For offline
preprocessing a centred window is possible because the whole clip is
available; the live pipeline uses an exponential moving average over past
frames instead.

\paragraph{6. Cropping and resizing.} The square region of side $L_t$ centred
on $c_t$ is cropped (with edge padding if it extends beyond the frame) and
resized to $96 \times 96$ pixels using area interpolation. The proposal allows
$96\times96$ or $112\times112$ \spec; $96\times96$ is chosen \rec\ because it
reduces input size by about 27\% and is comparable to the mouth-crop sizes used
in earlier sentence-level lipreading work on GRID~\cite{assael2016}. $112\times112$ remains an
option if the EfficientNet-B0 experiment benefits from it.

\paragraph{7. Normalising.} Crops are converted to grey-scale and stored as
8-bit integers. At load time they are scaled to $[0,1]$ and standardised with
the mean and standard deviation of all training crops. Grey-scale input
removes colour differences between speakers and cameras, which carry little
information about speech.

\paragraph{8. Forming sequences.} The crops of one utterance are stacked in
frame order into an array of shape \shape{T,96,96}. Temporal order is never
changed: data augmentation applies the same random crop offset and the same
horizontal flip to every frame of a clip, so the motion pattern is preserved.

\paragraph{9. Handling failures.} Three cases are handled explicitly.
\begin{itemize}
  \item \textbf{Occasional failed detections.} If Face Mesh fails on a frame,
        its box is linearly interpolated from the nearest successful frames on
        either side, and the frame is marked invalid in a per-frame mask.
  \item \textbf{Long runs of failures.} If more than 10\% of frames fail
        \rec, the utterance is excluded from training and listed in the
        integrity report. In the live demo, the visual stream for that window
        is treated as missing (\cref{sec:modality-dropout}).
  \item \textbf{Inconsistent lengths.} Clips are padded or truncated to the
        audio-derived frame count so that both streams describe the same time
        span; padding repeats the last valid frame and is marked in the mask.
\end{itemize}

\paragraph{Alignment with audio.} Because every crop retains its original frame
index and audio sample $n$ maps to frame $\lfloor 25n/16000\rfloor$, visual
frame $t$ is aligned with audio feature frames $4t$ to $4t+3$. This alignment
is preserved by never dropping frames: failed frames are interpolated, not
removed.

\subsection{Visual encoder architecture}
\label{sec:visual-encoder}
The proposal specifies a 3D convolutional front-end followed by a
bidirectional LSTM \spec, and its literature discussion also names
EfficientNetB0 as an efficient 2D backbone~\cite{kanavnet2026,ameer2025}. These
are two different ways of extracting per-frame spatial features, so this plan
defines them as two clearly separate configurations rather than stacking both
in one network.

\paragraph{Primary configuration V1: 3D-CNN + BiLSTM \rec.}
Three 3D convolution blocks process the clip. Each kernel spans several
consecutive frames, so every output feature responds both to the mouth's shape
and to how it changes over a short time window~\cite{lipsyncnet2024,assael2016}.
All temporal strides are 1, which keeps one feature vector per video frame;
only the spatial dimensions are downsampled. The spatial map of each frame is
then flattened, projected to 512 dimensions, and passed to a two-layer BiLSTM
that models longer-range context in both directions~\cite{kanavnet2026}. This
configuration is the proposal's ``3D convolutional visual encoder'' advanced
technique and is used in all main experiments.

\paragraph{Alternative configuration V2: 3D stem + EfficientNet-B0 trunk \tbd.}
A single 3D convolution captures short-term motion and outputs a 3-channel map
per frame. Each frame's map is then passed independently through an
ImageNet-pretrained EfficientNet-B0~\cite{tan2019} (by folding time into the
batch dimension), globally average-pooled to a 1280-dimensional vector,
projected to 512 dimensions and fed to the same BiLSTM. This mirrors the
common ``3D front-end plus 2D trunk'' design~\cite{afouras2022} and tests whether
transfer learning from a pretrained 2D backbone helps on GRID. Because
EfficientNet-B0 is applied to all 75 frames of every clip, V2 is several times
more expensive than V1; it is run as an ablation only if time and compute
allow (\cref{sec:schedule}).

\subsection{Tensor-shape walkthrough}
\Cref{tab:visual-shapes} traces one batch through configuration V1. The example
assumes a batch of $B=16$ clips, $T=75$ frames and $96 \times 96$ grey-scale
crops; the channel counts are recommended starting values \rec. Shapes follow
PyTorch's \shape{B,C,T,H,W} convention for 3D convolutions.

\begin{table}[H]
\centering
\caption{Visual encoder (configuration V1) tensor shapes for an example batch.}
\label{tab:visual-shapes}
\small
\begin{tabularx}{\textwidth}{@{}lXl@{}}
\toprule
\textbf{Stage} & \textbf{Operation} & \textbf{Output shape} \\
\midrule
Input & Normalised grey mouth crops & \shape{16,1,75,96,96} \\
Block 1 & Conv3d($1{\to}32$, kernel $3{\times}5{\times}5$, stride $1{\times}2{\times}2$), BN, ReLU & \shape{16,32,75,48,48} \\
        & MaxPool3d(kernel and stride $1{\times}2{\times}2$) & \shape{16,32,75,24,24} \\
Block 2 & Conv3d($32{\to}64$, kernel $3{\times}5{\times}5$, stride 1), BN, ReLU, MaxPool $1{\times}2{\times}2$ & \shape{16,64,75,12,12} \\
Block 3 & Conv3d($64{\to}96$, kernel $3{\times}3{\times}3$, stride 1), BN, ReLU, MaxPool $1{\times}2{\times}2$ & \shape{16,96,75,6,6} \\
Reshape & Move time forward and flatten each frame: $96\cdot6\cdot6 = 3456$ & \shape{16,75,3456} \\
Projection & Linear($3456{\to}512$), ReLU, Dropout(0.3) & \shape{16,75,512} \\
Temporal & 2-layer BiLSTM, 256 units per direction & \shape{16,75,512} \\
\bottomrule
\end{tabularx}
\end{table}

Temporal padding of 1 frame on each side of every convolution keeps the time
dimension at exactly 75, so the visual encoder emits one 512-dimensional vector
per video frame. The last dimension of the output, 512, is the concatenation of
the forward and backward LSTM states.

\subsection{Testing the visual pipeline}
\begin{itemize}
  \item A unit test runs the crop function on a stored GRID frame and asserts an
        output of shape $(96,96)$, type \texttt{uint8}, and a box centre inside the
        lower half of the face bounding box.
  \item A visual check script writes a contact sheet of crops for a random
        sample of utterances from every speaker; a team member inspects it for
        mis-centred or blurred crops before training starts.
  \item A unit test passes a random tensor of shape \shape{2,1,75,96,96} through
        the encoder and asserts an output of \shape{2,75,512} with finite values.
\end{itemize}
